\section{References}
\begingroup\small
\begin{itemize}[leftmargin=1.25em,itemsep=0.2em]
\item Sacred Scripture, Douay--Rheims/Challoner, American 1899 edition, eBible electronic witness: Matthew 5:43--48; 12:22--37; 24:45--51; 25; Luke 10:25--37; 12:35--48; 16:19--31; 23:39--43; John 5:21--29; Romans 2:1--24; 1 Corinthians 3:10--15; 4:1--7; Ephesians 2:1--10; James 2--3; 5:1--9; Revelation 20:11--15. Scriptural wording retains its public-domain status.
\item \emph{Catechism of the Catholic Church}, official English text: 678--679 (Christ as Judge); 1021--1022, 1030--1032, 1038--1041 (particular judgment, purification, and final judgment); 1451--1460, 1468--1473 (contrition, confession, satisfaction, reconciliation, and purification); 1735, 1790--1794 (responsibility and conscience); 2008--2011 (grace and merit); 2477, 2487--2489 (truth, reputation, reparation, and discretion).
\item Congregation for the Doctrine of the Faith, \href{https://www.vatican.va/roman_curia/congregations/cfaith/documents/rc_con_cfaith_doc_19790517_escatologia_en.html}{\emph{Letter on Certain Questions Concerning Eschatology}} (17 May 1979), nos.~2--5. Official English text; personal continuity and future bodily resurrection.
\item St.~Cyril of Jerusalem, \emph{Catechetical Lectures} XV.1--2, 23--26. Edwin Hamilton Gifford's translation, \emph{Nicene and Post-Nicene Fathers}, second series, vol.~7 (1894), retained CCEL volume transcription. The two comings, the impartial Judge, hidden fidelity, conscience, and works of mercy.
\item St.~Augustine, \emph{The City of God} XX.1--2, 14. Translation edited by Marcus Dods, vol.~2 (Edinburgh: T.~\& T.~Clark, 1871), retained Project Gutenberg eBook 45305 transcription. Final judgment, the present hiddenness of providence, and conscience before God.
\item St.~John Chrysostom, \emph{Homilies on Matthew} XLII.1--3, LXXVIII.1--3, and LXXIX.1--2. George Prevost's translation, revised by M.~B.~Riddle, \emph{Nicene and Post-Nicene Fathers}, first series, vol.~10 (1888), retained CCEL transcription. Accountable speech and examination; the virgins, talents, and works of mercy. The online volume's damaged transition into homily LXXIX is excluded from use; the cited material from its first section is the intact later discussion of accessible works and poor believers.
\item St.~Thomas Aquinas, \emph{Summa theologiae} III, q.~59, aa.~2--3, 5; I, q.~21, a.~4. Fathers of the English Dominican Province translation, Project Gutenberg eBooks 19950 and 17611 (2006 electronic witnesses). Christ's judicial authority, continuing historical consequences without new postmortem merit, and mercy at the root of divine justice.
\item \emph{Summa theologiae}, Supplement, q.~87, a.~2, body and reply~3. \emph{Summa Theologica, Third Part (Supplement), Questions LXXXVII--XCIX and Appendices}, English Dominican translation (London: Burns Oates and Washbourne; New York: Benziger Brothers, 1922), pp.~5--7. Posthumous compilation drawn from Thomas's earlier teaching, principally the commentary on the \emph{Sentences}. Disclosure of forgiven sins is presented as the more probable opinion, with Peter Lombard's contrary position expressly recorded.
\item St.~Francis de Sales, \emph{Introduction to the Devout Life} I.14, especially considerations~5--6, affections and resolutions~2, and conclusion. \emph{A New Translation} (London, Oxford, and Cambridge: Rivingtons, 1876), pp.~38--40; translator unnamed in this printing. Conscience, the two sentences, confession, and gratitude for time to repent.
\item St.~Alphonsus Liguori, \emph{Preparation for Death; or, Considerations on the Eternal Maxims}, consideration XXIV, point~3, especially p.~299. Translated by ``a Catholic clergyman'' (Dublin: James Duffy, 1843), pp.~296--300. Examination and conversion before judgment.
\item Benedict XVI, \href{https://www.vatican.va/content/benedict-xvi/en/encyclicals/documents/hf_ben-xvi_enc_20071130_spe-salvi.html}{\emph{Spe salvi}} (2007), nos.~47--48. Official English text: the theological proposal of purifying encounter and the communion expressed in intercession for the dead.
\end{itemize}\endgroup
